\section{Related Work}
\label{sec:related}
% About 0.75 page. Each paper is described only from its abstract (arXiv or ACL Anthology).

\paragraph{Chunking for retrieval.} Several studies compare many
chunking methods. \citet{smigielski2026chunking} evaluate a wide range of
methods against their computational cost, and \citet{shaukat2026systematic}
benchmark fixed-size, semantic, structure-aware, hierarchical, adaptive and
LLM-assisted methods across domains and embedding models, finding that
content-aware chunking beats naive fixed-length splitting.
\citet{amiri2025chunk} study how chunking interacts with the embedding
model for chemistry papers. Other work builds new chunkers:
\citet{zhao2025moc} measure chunk quality by boundary clarity and chunk
stickiness and train a mixture of LLM chunkers, and \citet{lu2025hichunk}
argue that sparse evidence makes existing benchmarks weak for chunking and
propose a benchmark with multi-level chunk points and a hierarchical
chunker. Most of these comparisons fix the number of retrieved chunks; we
compare at an equal token budget, where chunk size and retrieval depth
trade off.

\paragraph{Choosing per document or per task.}
\citet{demourajunior2026adaptive} select a chunking method per document
from intrinsic, document-based quality metrics; we reuse their recursive
separators and ask whether per-document choice beats one fixed setup when
both are scored on held-out questions. \citet{orbach2025hpo} benchmark
hyper-parameter optimisation for retrieval-augmented generation and find
that greedy or random search over the configuration is efficient and
improves performance on every dataset they test.

\paragraph{Parsing, OCR and document structure.} Docling
\citep{auer2024docling} converts PDFs into structured documents;
OmniDocBench \citep{ouyang2024omnidocbench} benchmarks such parsers on
diverse document types. OHR-Bench \citep{zhang2024ocr} measures how OCR
errors cascade into retrieval-augmented generation, and
\citet{sun2026ocr} find that high OCR accuracy does not guarantee strong
downstream retrieval. \citet{santos2026pdf} compare PDF-to-Markdown
frameworks on Portuguese administrative documents, and
\citet{elbachyr2026empirical} study parsing and chunking choices for
financial question answering. \citet{lin2026pdf} retrieve structure-aware
evidence units from clinical guidelines, and \citet{singh2026parser}
evaluate parsers, chunking strategies and embedding models jointly on
Indian regulatory documents, whose questions we use. Our study differs in
three ways: the analysis plan, scorer and decision rule were frozen before
the held-out documents were opened, the per-document oracle is
cross-validated within each document, and every comparison uses the same
token budget.
